% GENERATED FILE. Edit canonical lessons, then run npm run generate:books.
% canonical-source-hash: fnv1a64:a2065ccc5cb97db8
% canonical-lessons: GU-C476-ghorvu, GU-C476-nidar, GU-C476-jagran, GU-C476-podhvu, GU-C476-anidra

\chapter{A Long Night's Sleep}
\label{ch:gu-a-long-night-s-sleep}
\hlchaptermodality{476}

\begin{chapteropening}
\textbf{By the end of this chapter:} \emph{I can talk about how well I slept.}

Complete the last lesson of chapter 476: I can talk about how well I slept.
\end{chapteropening}

\section[\texorpdfstring{\gu{ઘોરવું} (ghorvũ)}{ઘોરવું (ghorvũ)}]{\gu{ઘોરવું} (ghorvũ) — to snore}

\label{lesson:GU-C476-ghorvu}



\begin{quote}
Before the new one: say the Gujarati for eyesight, then the Gujarati for to touch.
\end{quote}

\subsection*{You'll want to know: \gu{ઘોરવું}}
\textbf{\gu{ઘોરવું}} — \emph{ghorvũ} — \textquotedblleft{}to snore\textquotedblright{}.

\begin{quote}
\gu{દાદા} \gu{આખી} \gu{રાત} \gu{ઘોરે} \gu{છે}.
\end{quote}

\begin{quote}
\emph{dādā ākhī rāt ghore che.}
\end{quote}

\begin{quote}
Grandfather snores all night.
\end{quote}

Past (he …): \gu{તે} \gu{આખી} \gu{રાત} \gu{ઘોર્યો}, \emph{te ākhī rāt ghoryo}.

\begin{cousinweb}
\textbf{\gu{ઘોરવું}} (\emph{ghorvũ}) goes back to Sanskrit \emph{ghurati}, \textquotedblleft{}cries out frighteningly\textquotedblright{}, probably an imitative word; CDIAL lists the Gujarati verb, \textquotedblleft{}to snore\textquotedblright{} (CDIAL 4487).
\end{cousinweb}

\begin{grammarlens}[title={What you've built}]
One more word for describing what you lived and felt.
\end{grammarlens}

\subsection*{Your turn}
\begin{itemize}
\raggedright
  \item \emph{Say it:} \emph{ghorvũ}
  \item \emph{Say it:} \emph{ghorvũ}, once more
  \item \emph{Say it:} \emph{dādā ākhī rāt ghore che.} — Grandfather snores all night.
  \item \emph{Recall:} say the Gujarati for eyesight, then the Gujarati for to touch, then say \emph{ghorvũ} again
\end{itemize}

\begin{tcolorbox}[breakable,colback=teal!4,colframe=teal!35!black,title={Before you move on}]
What does \gu{ઘોરવું} mean? (\textquotedblleft{}To snore\textquotedblright{}.) Say it once more.
\end{tcolorbox}
\section[\texorpdfstring{\gu{નીંદર} (nī̃dar)}{નીંદર (nī̃dar)}]{\gu{નીંદર} (nī̃dar) — slumber; sleep (a homely word)}

\label{lesson:GU-C476-nidar}



\begin{quote}
Before the new one: say the Gujarati for to touch, then the Gujarati for to snore.
\end{quote}

\subsection*{You'll want to know: \gu{નીંદર}}
\textbf{\gu{નીંદર}} — \emph{nī̃dar} — \textquotedblleft{}slumber; sleep (a homely word)\textquotedblright{}.

\begin{quote}
\gu{થાકને} \gu{લીધે} \gu{મને} \gu{તરત} \gu{નીંદર} \gu{આવી} \gu{ગઈ}.
\end{quote}

\begin{quote}
\emph{thākne līdhe mane tarat nī̃dar āvī gaī.}
\end{quote}

\begin{quote}
I was so tired that I fell asleep at once.
\end{quote}

A feminine noun.

\begin{cousinweb}
\textbf{\gu{નીંદર}} (\emph{nī̃dar}) is inherited from Sanskrit \emph{nidrā}, \textquotedblleft{}sleep\textquotedblright{}; CDIAL lists the Gujarati \emph{nīdar} (CDIAL 7200).
\end{cousinweb}

\begin{grammarlens}[title={What you've built}]
One more word for describing what you lived and felt.
\end{grammarlens}

\subsection*{Your turn}
\begin{itemize}
\raggedright
  \item \emph{Say it:} \emph{nī̃dar}
  \item \emph{Say it:} \emph{nī̃dar}, once more
  \item \emph{Say it:} \emph{thākne līdhe mane tarat nī̃dar āvī gaī.} — I was so tired that I fell asleep at once.
  \item \emph{Recall:} say the Gujarati for to touch, then the Gujarati for to snore, then say \emph{nī̃dar} again
\end{itemize}

\begin{tcolorbox}[breakable,colback=teal!4,colframe=teal!35!black,title={Before you move on}]
What does \gu{નીંદર} mean? (\textquotedblleft{}Slumber; sleep (a homely word)\textquotedblright{}.) Say it once more.
\end{tcolorbox}
\section[\texorpdfstring{\gu{જાગરણ} (jāgraṇ)}{જાગરણ (jāgraṇ)}]{\gu{જાગરણ} (jāgraṇ) — staying up all night}

\label{lesson:GU-C476-jagran}



\begin{quote}
Before the new one: say the Gujarati for to snore, then the Gujarati for slumber.
\end{quote}

\subsection*{You'll want to know: \gu{જાગરણ}}
\textbf{\gu{જાગરણ}} — \emph{jāgraṇ} — \textquotedblleft{}staying up all night\textquotedblright{}.

\begin{quote}
\gu{દર્દીને} \gu{સંભાળવા} \gu{મેં} \gu{આખી} \gu{રાત} \gu{જાગરણ} \gu{કર્યું}.
\end{quote}

\begin{quote}
\emph{dardīne sambhāḷvā mẽ ākhī rāt jāgraṇ karyũ.}
\end{quote}

\begin{quote}
I stayed up all night to look after the patient.
\end{quote}

A neuter noun.

\begin{cousinweb}
\textbf{\gu{જાગરણ}} (\emph{jāgraṇ}) is a tatsama, Sanskrit \emph{jāgaraṇa}, \textquotedblleft{}waking, keeping awake\textquotedblright{}, as Monier-Williams records it.
\end{cousinweb}

\begin{grammarlens}[title={What you've built}]
One more word for describing what you lived and felt.
\end{grammarlens}

\subsection*{Your turn}
\begin{itemize}
\raggedright
  \item \emph{Say it:} \emph{jāgraṇ}
  \item \emph{Say it:} \emph{jāgraṇ}, once more
  \item \emph{Say it:} \emph{dardīne sambhāḷvā mẽ ākhī rāt jāgraṇ karyũ.} — I stayed up all night to look after the patient.
  \item \emph{Recall:} say the Gujarati for to snore, then the Gujarati for slumber, then say \emph{jāgraṇ} again
\end{itemize}

\begin{tcolorbox}[breakable,colback=teal!4,colframe=teal!35!black,title={Before you move on}]
What does \gu{જાગરણ} mean? (\textquotedblleft{}Staying up all night\textquotedblright{}.) Say it once more.
\end{tcolorbox}
\section[\texorpdfstring{\gu{પોઢવું} (poḍhvũ)}{પોઢવું (poḍhvũ)}]{\gu{પોઢવું} (poḍhvũ) — to lie down to sleep}

\label{lesson:GU-C476-podhvu}



\begin{quote}
Before the new one: say the Gujarati for slumber, then the Gujarati for staying up all night.
\end{quote}

\subsection*{You'll want to know: \gu{પોઢવું}}
\textbf{\gu{પોઢવું}} — \emph{poḍhvũ} — \textquotedblleft{}to lie down to sleep\textquotedblright{}.

\begin{quote}
\gu{બાળક} \gu{પલંગ} \gu{પર} \gu{પોઢી} \gu{ગયું}.
\end{quote}

\begin{quote}
\emph{bāḷak palaṅg par poḍhī gayũ.}
\end{quote}

\begin{quote}
The baby fell asleep on the bed.
\end{quote}

Past (it …): \gu{તે} \gu{બાળક} \gu{પોઢી} \gu{ગયું}, \emph{te bāḷak poḍhī gayũ}.

\begin{cousinweb}
\textbf{\gu{પોઢવું}} (\emph{poḍhvũ}) is inherited: CDIAL lists the Gujarati verb, \textquotedblleft{}to sleep\textquotedblright{}, under Sanskrit \emph{pravardhate}, \textquotedblleft{}grows up\textquotedblright{}, whose sense \textquotedblleft{}lie down\textquotedblright{} CDIAL explains as \textquotedblleft{}stretched out\textquotedblright{} (CDIAL 8789).
\end{cousinweb}

\begin{grammarlens}[title={What you've built}]
One more word for describing what you lived and felt.
\end{grammarlens}

\subsection*{Your turn}
\begin{itemize}
\raggedright
  \item \emph{Say it:} \emph{poḍhvũ}
  \item \emph{Say it:} \emph{poḍhvũ}, once more
  \item \emph{Say it:} \emph{bāḷak palaṅg par poḍhī gayũ.} — The baby fell asleep on the bed.
  \item \emph{Recall:} say the Gujarati for slumber, then the Gujarati for staying up all night, then say \emph{poḍhvũ} again
\end{itemize}

\begin{tcolorbox}[breakable,colback=teal!4,colframe=teal!35!black,title={Before you move on}]
What does \gu{પોઢવું} mean? (\textquotedblleft{}To lie down to sleep\textquotedblright{}.) Say it once more.
\end{tcolorbox}
\section[\texorpdfstring{\gu{અનિદ્રા} (anidrā)}{અનિદ્રા (anidrā)}]{\gu{અનિદ્રા} (anidrā) — sleeplessness}

\label{lesson:GU-C476-anidra}



\begin{quote}
Before the new one: say the Gujarati for staying up all night, then the Gujarati for to lie down to sleep.
\end{quote}

\subsection*{You'll want to know: \gu{અનિદ્રા}}
\textbf{\gu{અનિદ્રા}} — \emph{anidrā} — \textquotedblleft{}sleeplessness\textquotedblright{}.

\begin{quote}
\gu{ચિંતાને} \gu{લીધે} \gu{મામાને} \gu{અનિદ્રાનો} \gu{રોગ} \gu{છે}.
\end{quote}

\begin{quote}
\emph{cintāne līdhe māmāne anidrāno rog che.}
\end{quote}

\begin{quote}
Because of worry my uncle suffers from sleeplessness.
\end{quote}

A feminine noun.

\begin{cousinweb}
\textbf{\gu{અનિદ્રા}} (\emph{anidrā}) is a tatsama, Sanskrit \emph{anidrā}, \textquotedblleft{}sleeplessness\textquotedblright{}, as Monier-Williams records it: \emph{a-}, \textquotedblleft{}not\textquotedblright{}, and \emph{nidrā}, \textquotedblleft{}sleep\textquotedblright{}.
\end{cousinweb}

\begin{grammarlens}[title={What you've built}]
One more word for describing what you lived and felt.
\end{grammarlens}

\subsection*{Your turn}
\begin{itemize}
\raggedright
  \item \emph{Say it:} \emph{anidrā}
  \item \emph{Say it:} \emph{anidrā}, once more
  \item \emph{Say it:} \emph{cintāne līdhe māmāne anidrāno rog che.} — Because of worry my uncle suffers from sleeplessness.
  \item \emph{Recall:} say the Gujarati for staying up all night, then the Gujarati for to lie down to sleep, then say \emph{anidrā} again
\end{itemize}

\begin{tcolorbox}[breakable,colback=teal!4,colframe=teal!35!black,title={Before you move on}]
What does \gu{અનિદ્રા} mean? (\textquotedblleft{}Sleeplessness\textquotedblright{}.) Say it once more.
\end{tcolorbox}
